\documentclass[twoside]{article}

\usepackage{ustj}
\usepackage{enumitem}

\addbibresource{mss.bib}

\newcommand{\authorname}{Evan Forman~\patp{bonbud-macryg} \\
Jake Hamilton~\patp{tondes-sitrym} \\
Trent Steen~\patp{hanfel-dovned} \\
}
\newcommand{\affiliation}{Groundwire Foundation}
\newcommand{\authorpatp}{\patp{bonbud-macryg}}

%  Make first page footer:
\fancypagestyle{firststyle}{%
\fancyhf{}% Clear header/footer
\fancyhead{}
\fancyfoot[L]{{\footnotesize
              %% We toggle between these:
              % Manuscript submitted for review.\\
              {\it Urbit Systems Technical Journal} III:2 (2026):  87–122. \\
              ~ \\
              Address author correspondence to \authorpatp.
              }}
}
%  Arrange subsequent pages:
\fancyhf{}
\fancyhead[LE]{{\urbitfont Urbit Systems Technical Journal}}
\fancyhead[RO]{Groundwire}
\fancyfoot[LE,RO]{\thepage}

%%MANUSCRIPT
\title{Groundwire: Decentralized Sybil Resistance for a Dead Internet}
\author{\authorname \\ \affiliation}
\date{}

\begin{document}

\maketitle
\thispagestyle{firststyle}

\begin{abstract}
\sloppy
Public-key infrastructures have hitherto assured us that a message comes
from whoever it says it does, but we contend that this does not give us
secure, sybil-resistant networking at the current level of \textsc{llm}
capability and proliferation, let alone post- or even near-\textsc{agi}
conditions. Once bots represent a majority of internet traffic we have
to ask, who is trying to send us this message? As \textsc{llm}s become capable of
sophisticated social engineering attacks we must ask, is the sender of
this message really who they say they are? And as AI drives the price of
information production to near-zero, we ask: at what cost? This article
introduces a decentralized, sybil-resistant identity and peer-to-peer
networking stack which we design for an internet where the overwhelming
majority of participants are autonomous AI agents.
\end{abstract}

% We will adjust page numbering in final editing.
\pagenumbering{arabic}
\setcounter{page}{87}

\tableofcontents

\section{Introduction}\label{introduction}

Spam prevention efforts by legacy web platforms have failed to prevent the popular adoption of ``dead internet theory''---that most visible
internet activity is from bots---and an incipient retreat from the
public web into closed (if not exactly private) online spaces.
Everywhere you go online, there you are: ticking a box to prove you are
human. Public discourse and culture at large suffer as the information
environment is increasingly beset by engagement-farming bot accounts
and increasingly automated, agentic pipelines for state and
non-state manipulation of target demographics \citep{Lumbaca2025}.

Spam prevention on these platforms is stymied by fundamental conflicts
of interest: as AI agents become more capable, it becomes harder to
automatically discriminate between human and nonhuman input, so human
users suffer more onerous verification burdens. This runs counter to the
platforms' incentive to attract and retain as many users for as much
time-spend as possible, as well as their incentive to report as many
impressions to creators and advertisers as possible.

How these platforms would adapt to an internet where an increasing
supermajority of participants are AI agents is anyone's guess: \textsc{sim}
verification and photo ID laws such as those in Japan and Australia just
increase operating expenses in the seller's market for fake accounts.
Fake accounts tied to phone numbers in those markets currently sell for
an average \$4.93 (Japan) and \$3.24 (Australia) per account in the same
bulk volumes as cheaper \$0.10 (UK) and \$0.26 (US) accounts, as sellers
pass \textsc{sim} costs to buyers, with prices for accounts on Telegram and WhatsApp rising
up to 15\% in the lead-up to national elections \citep{Cambridge2025}.

The spam resistance stack we outline here is general purpose: any
platform could implement some or all of it. Our interest lies not in
redeeming Web 2.0 from its malign incentives, nor do we seek a return
to a prelapsarian ``open web'' whose very openness exposed it to
predation, centralization, and enclosure \citep{Davis2022}.

This article introduces three separable but complementary layers of a new
decentralized identity and peer-to-peer networking stack:

\begin{enumerate}[noitemsep,topsep=2pt]
\def\labelenumi{\arabic{enumi}.}
\item
  \sloppy
  Mnemonyms \S\ref{mnemonyms}, a \textsc{bip}-39-style (\citeyear{BIP0039}) mnemonic system
  for reasonably short, human-meaningful references to public keys.
\item
  \sloppy
  Groundwire ID \S\ref{groundwire-id}, a protocol for unique, persistent cryptographic
  identity across key rotations, tied to public signatures.
\item
  \sloppy
  Wire \S\ref{wire-uris}, a \textsc{uri} scheme for signing data on arbitrary namespaces,
  securing it against the value of the signer's Groundwire ID.
\end{enumerate}

In the Future Work section \S\ref{future-work} we outline near- to
medium-term projects including named data networking secured against
Groundwire ID, post-quantum signatures, and derivative identities for
AI agents and networked devices.

Appendices \S\ref{appendices} detail our thinking around two key
issues for the software ecosystem we anticipate: security by planned
ossification, and trustless consensus mechanisms.

\sloppy
Groundwire is in part an argument about structure: we contend that the
collective cultivation of a global filesystem-like (or, ``tree-shaped'')
namespace where end-users are cryptographically sovereign root nodes can
deliver a more effective sybil resistance regime than all mainstream
solutions---\textsc{captcha}, \textsc{sim}/ID/Email verification,
blacklisting IP ranges, etc.---which are isolated, interactive,
``conversation-'' or ``connection-oriented'' exchanges, as opposed to
our ``data-oriented'' (as in Named Data Networking,
\citeauthor{DirectedMessaging2026}, 2026) approach to sybil resistance.

\sloppy
Materially, Groundwire's first implementation \S\ref{groundwire-id}
requires a one-time proof-of-work action from the creator of a
self-custodied keypair to attest to owning that keypair on the Bitcoin
blockchain. This sybil resistance ``challenge'' only happens once, and
is cashed out in every network interaction that involves the user
sending an effectful write or an access-controlled read
\S\ref{named-data-networking}. Server operators may layer further
automated checks on top of this, but no manual sybil resistance checks
should be necessary to prove that the end-user is not a disposable
botnet identity. If servers want to verify that the user is a human and
not a benign AI agent, that's their business; the protocol does not
discriminate.

\sloppy
A client may reconcile the namespace from multiple sources if it is not
already unified in-situ. In the sybil resistance case they may apply
policy to logical subtrees, analyse node provenance across space (permissions,
derivative agent identities \S\ref{hierarchical-identity}) and time
(key rotations, verifiable onchain history \S\ref{onchain-attestation}). When this exact same
filesystem addresses data storage, \mbox{financial} transactions, and
networking traffic \S\ref{named-data-networking}, the design space for
simple and efficient (so, decentralizable) sybil resistance tooling is wide open.
Any subtree in the namespace can be summed into some quantity of trust-value over and
above the total financial, computational, and/or thermodynamic cost of
creating the nodes in that subtree, where every node in the subtree is
ultimately secured against the proof-of-work value of the keypair that
created it.

We take care from the outset to support interoperability between
software that uses Mnemonyms, Groundwire ID, and Wire \textsc{uri}s,
preferring to leave composition of that software to developers and end-users.

The sections below are introductory snapshots of works in progress, and
not full specifications. The Groundwire Foundation has implemented these
protocols in public alpha software, and everything described below is
subject to change before production launch.

\section{Mnemonyms}\label{mnemonyms}

The mnemonym scheme is a general-purpose method for generating mnemonic
references to binary strings. We use these as reasonably short,
human-meaningful references to public keys on a \textsc{pki} such as Bitcoin,
Nostr, and Groundwire ID.

The mnemonym design is greatly inspired by \textsc{bip}-39. Like
\textsc{bip}-39, a mnemonym wordlist is an alphabetically-sorted list
of 2,048 words. Each word in the string represents 11 bits in a
bitstring of entropy $\textsc{ent}$ with an appended checksum of
bitlength $(\textsc{ent}/32)$.

Two key differences from \textsc{bip}-39 mnemonics:

\begin{itemize}[noitemsep,topsep=2pt]

\item
  We do not expect mnemonym wordlists to be embedded in crypto hardware
  wallets in the base case; as a result, we grant a proliferation of
  wordlists for various languages, communities, and personal
  use-cases.
\item
  Mnemonym strings represent an entire series of leading null words by
  omission; non-leading null words are represented normally, with
  whichever word is at index 0 in the list.
\end{itemize}

By omitting leading zeroes, we allow for reasonably short mnemonyms (or
``nyms'') to be distributed to users according to whichever scheme the
underlying \textsc{pki} permits. Such schemes may include simple
first-come-first-serve domain registration, digital
feudalism \citep{YarvinSpecIntro}, or proof-of-work
mining. In a centralized, permissioned, or feudal system the length
of the nym is useless as a market signal in and of itself: you must
understand its provenance to find out how it was attained and at what cost,
if any. In a permissionless proof-of-work system, the length of the nym
approximately illustrates the thermodynamic cost at which the keypair
must have been discovered and acquired, whether the holder of the nym
mined it on their own machine or purchased it from a miner.

Our initial English wordlist \S\ref{wordlist-construction} consists
exclusively of iambs. For example, the pubkey\ldots{}

\begin{lstlisting}[style=listingcode]
00000000000ba839b38bcfd68037176c
\end{lstlisting}

\noindent
would be encoded as\ldots{}

\begin{lstlisting}[style=listingcode]
.preserves.astounds.regrets.ensoul.repeats.about
.receive.repulsed
\end{lstlisting}

\noindent
We reserve five visual representations of nyms at the protocol level:

\begin{itemize}[noitemsep,topsep=2pt]

\item
  Bare nyms, e.\,g.,
  \par{\centering
  \texttt{abet.baboon.caffeine.denounce.escape}\par}
\item
  One-dot nyms, e.\,g.,
  \par{\centering
  \texttt{.abet.baboon.caffeine.denounce.escape}\par}
\item
  Two-dot nyms, e.\,g.,
  \par{\centering
  \texttt{..abet.baboon.caffeine.denounce.escape}\par}
\item
  ``Foreshortened'' nyms of a string's first, second, penultimate, and
  final words, e.\,g.,
  \par{\centering
  \texttt{\phantom{~}.abet.baboon..denounce.escape}\\
  \texttt{..abet.baboon..denounce.escape}\par}
\item
  ``Abridged'' nyms of just its first and last words,
  e.\,g.,
  \par{\centering
    \texttt{\phantom{~}.abet...escape}\\
    \texttt{..abet...escape}\par}
\end{itemize}

Applications must reserve these as follows:

\begin{itemize}[noitemsep,topsep=2pt]

\item
  One-dot for nyms which they have verified as being on the relevant
  \textsc{pki}.
\item
  Two-dot for nyms which they have either confirmed are not on the \textsc{pki},
  or for disposable references they do not expect to correspond to user
  identities on the \textsc{pki} (e.\,g.\ block hashes, content hashes, receiving
  addresses, etc.).
\item
  Bare nyms are reserved for intermediate contexts where a sender and
    receiver may or may not respect the same \textsc{pki} (e.\,g.\ Wire \textsc{uri}s \S\ref{wire-uris}).
\end{itemize}

``Foreshortened'' and ``abridged'' nyms may be used however the
application developer sees fit, though abridged nyms must not be used in
contexts where user funds may be at risk, to mitigate against address poisoning.

These are the only nym representations we formally bless. Informally, we
grant that pseudonymous users may come up with various colloquial,
context-sensitive shortname conventions to refer to themselves and
eachother, such as ``\texttt{.abet.escape}'', ``\texttt{caffeine.chanteuse}'',
``\texttt{.eclipse}'', ``\texttt{eclipse}'', etc. We expect developers will layer
conventional ``contact book'' or ``display name'' features on top of
mnemonyms at some point, obviating their use in online social contexts
beyond first contact or peer-to-peer financial transactions.

The iambic English wordlist and libraries in C, Python, and Hoon are available at
\href{https://github.com/gwbtc/mnemonyms}{\texttt{gwbtc/mnemonyms}}. Our C library is implemented in
\href{https://github.com/gwbtc/comet-miner}{\texttt{gwbtc/comet-miner}}. The
iambic English wordlist is still a work-in-progress and likely to change.

\subsection{Multi-line layouts}\label{multi-line-layouts}

Application developers may lay the nym out across several lines, e.\,g.\ in
a UI context specifically designed to help the user memorize a nym. We
encourage the following layouts specifically:

Twelve-word nym on three lines of four words:

\begin{lstlisting}[style=listingcode]
.eclipse.betray.escape.delight
.mankind.alone.delay.ignite
.implode.elect.enclose.alas
\end{lstlisting}

Ten-word nym on two lines of five words:

\begin{lstlisting}[style=listingcode]
.tattoo.giraffe.kazoo.champagne.balloon
.auteurs.ablaze.admire.untamed.display
\end{lstlisting}

Nine-word nym on three lines of three words:

\begin{lstlisting}[style=listingcode]
.raccoon.beret.lagoon
.typhoon.guitar.chateau
.concur.exempt.vacate
\end{lstlisting}

Eight-word nym on two lines of four words:

\begin{lstlisting}[style=listingcode]
.disgrace.expired.infest.reject
.doubloon.retells.entice.retreat
\end{lstlisting}

\subsection{Wordlist construction}\label{wordlist-construction}

A nym wordlist may conform to any particular aesthetic or none at all.
Wordlists should aspire to make resultant nyms as memorable, legible,
and easy to work with as possible.

In our iambic English wordlist, we achieve this by trying to fulfill the
following criteria:

\begin{itemize}[noitemsep,topsep=2pt]

\item
  Every word is a two-syllable iamb, which enhances memorability in that
  it reduces the number of words that could possibly follow another
  without reference to the wordlist.
\item
  Words have no punctuation or diacritics in them.
\item
  We limit the number of words that share common four-letter prefixes
  like ``over-'', ``semi-'', and ``tran-'' for the sake of
  autocomplete features in text inputs.
\item
  No two words are exact homophones.
\item
  We err on the side of words which are relatively common.
\item
  Many words are plurals or verbs that end with `s', so we limit words
  beginning with `s' to avoid confusion in speech.
\item
  We err on the side of words which result in absurd, concretely visual
  combinations inspired by memory palace techniques, though it should be
  just as easy to come up with a nym that is beautiful as one that is
  humorous.
\end{itemize}

\section{Groundwire ID}\label{groundwire-id}

Groundwire ID is a platform-agnostic protocol for persistent
decentralized identities belonging to humans and autonomous AI agents.

The user is assumed to hold a \textsc{bip}-32 (\citeyear{BIP0032}) HD crypto wallet
that can derive an Ed25519 keypair whose pubkey will be encoded as a
mnemonym. (Groundwire Foundation software uses the iambic English
wordlist, but no wordlist is specified at the protocol level.
Applications may use whichever wordlist they choose or let the user
configure one.)

In practice, Groundwire ID is a protocol for tweaking pubkeys with
publicly-verifiable evidence of their sybil resistance value, where we
take the position that ``sybil resistance value'' is mostly
financial/thermodynamic cost plus supplementary analytics
subject to some policy configured by the user that determines who can send
packets their way \S\ref{offchain-revelation}.

This tweak has two components: \textsc{pki} metadata, and a pointer to
publicly-verifiable evidence on that \textsc{pki}. As we outlined in
\textsc{uip}-0136 (\citeyear{UIP-0136}), the tweak is a
``bidirectional cryptographic commitment between onchain and offchain
attestations to prevent replay, spoofing, and double-spawning attacks. If
there wasn't an offchain anchor pointing to onchain data, then anyone
could come in and attest to ownership of any comet and \texttt{\%ord-watcher}
would naively believe them.''

\begin{lstlisting}[style=listingcode]
// pseudocode to produce a Groundwire ID tweak
(concatenate [(length-encode 'urbtc')
                [(length-encode 8)
                 (length-encode
                  (serialize pki-evidence))]])
\end{lstlisting}

Where \texttt{(concatenate\ ...)} concatenates an array of
\texttt{{[}bitwidth\ bitstring{]}} pairs. This produces a
bitstring which is hashed into the pubkey to couple it to the \textsc{pki} evidence.

\begin{lstlisting}[style=listingcode]
// pseudocode to apply the tweak to a pubkey
(sha256 pubkey
  (concatenate [(length-encode 'urbtc')
                  [(length-encode 8)
                   (length-encode
                    (serialize pki-evidence))]]))
\end{lstlisting}

The first two items in this array are the \textsc{pki} metadata \S\ref{pki-metadata}, and the
third a link to \textsc{pki} evidence \S\ref{pki-evidence}.

\subsection{Personally identifiable
information}\label{personally-identifiable-information}

Notably, Groundwire ID does not imply that a verified nym correlates to
a human person. (We assume an internet where 99.9\% of users are
autonomous AI agents, an addressable market easily served by the 256-bit
namespace that Ed25519 nyms imply and plausibly served by a 128-bit
namespace \S\ref{hierarchical-identity}.)

It \emph{could} verifiably correlate to a human, and we are interested
in exploring permissionless ZK solutions to couple Groundwire IDs to personally identifiable information
without exposing that \textsc{pii} to anyone. A bridge towards \textsc{pii} in the general
case and government-backed ID in particular is one direction in a
greater design space: What can we cryptographically attest about an
identity? And how can we compose those predicates into subprotocols
useful for a variety of autonomous network actors within the bounds of a
simple, freezable format?

\subsection[\textsc{pki} metadata]{PKI metadata}\label{pki-metadata}

A \textsc{pki} may be one of several on the same technical substrate, e.\,g.\ the
same blockchain. The tweak names this \textsc{pki} by a string of
\texttt{\^{}{[}a-z{]}{[}a-z0-9-{]}*\$}. In the
\texttt{(length-encode\ \textquotesingle{}urbtc\textquotesingle{})}
example above, the name of the \textsc{pki} is ``urbtc'', which is Groundwire's
Urbit \textsc{pki} on Bitcoin. Verifiers will resolve this name however
appropriate: in urbtc, the \textsc{pki} name is also the name of the local
application an Urbit node will use to verify identities on that \textsc{pki}, which may
be one of several \textsc{pki}s they follow.

The other piece of \textsc{pki} metadata included in the tweak is
\texttt{(length-encode\ 8)}, which in this case tells us that the pubkey
was tweaked against urbtc version 8K. \textsc{pki}s may support some kelvin
versions \S\ref{ossification} and not others, and may have their own control flow for
handling tweaks on old versions.

One purpose of this is that it namespaces Groundwire IDs across \textsc{pki}s:
the same inputs will not generate the same Groundwire ID on two
different \textsc{pki}s, even if they're on the same substrate. The probability
of the same Groundwire ID (that is, the same verified nym) appearing on two
\textsc{pki}s is astronomically tiny. This improbability leaves open a design
space for cross-\textsc{pki} interoperability, e.\,g.\ ``teleburning'' Ethereum
assets onto Bitcoin Ordinal inscriptions \citep{Rodarmor2026Teleburning}, but at time of
writing we have no plans to pursue this.

\subsection[\textsc{pki} evidence]{PKI evidence}\label{pki-evidence}

The second part of the tweak is some publicly-verifiable data on the \textsc{pki}
for a verifier to do two things: 1) verify the existence of a peer on
the \textsc{pki} and 2) judge the sybil resistance value of that data, according
to some policy hard-coded in the verifier or set by the user.

The content of the \texttt{pki-evidence} in the
\texttt{(length-encode\ (serialize\ pki-evidence))} pseudocode above will vary by PKI.
For example, in urbtc \S\ref{urbtc-implementation} the \textsc{pki} evidence is the
relative ordinal number \citep{Rodarmor2022Ordinals} of a sat
in a \textsc{utxo} belonging to the Bitcoin wallet that generated the Groundwire
ID in question. We avoid using absolute ordinal numbers for performance
reasons.

Notably, we do not include key rotations here or anywhere in the tweak,
because doing so would change the resulting nym every key rotation.
Reconciling the history of an identity across key rotations is the job
of the \textsc{pki}. Key rotations should generally not be considered to
``reset'' the sybil resistance value of an ID: on a blockchain where
operations require transaction fees paid to infrastructure providers, an
ID should accrue sybil resistance value over successive key rotations
and miscellaneous transactions.

\subsection{urbtc implementation}\label{urbtc-implementation}

\sloppy
The centerpiece of our public alpha is urbtc, a protocol for maintaining
an Urbit \textsc{pki} on the Bitcoin blockchain. urbtc only supports Urbit
``comets''---nodes in Urbit's self-signed 128-bit namespace---which
can be spawned by anyone and have no speculative value. Ordinarily,
comets are reliant on one of \mbox{Urbit's} 256 ``galaxy'' root nodes to do
peer discovery and routing, and the comet's galaxy is hard-coded by the
last 8 bits of their ID. Groundwire allows any comet to serve as a root
node by attesting to a stable IP address onchain.

An urbtc tweak looks like this:

\begin{lstlisting}[style=listingcode]
// pseudocode
(concatenate [(length-encode 'urbtc')
                [(length-encode 8)
                 (length-encode
                  (serialize spawn-sont))]])
\end{lstlisting}

Where the \texttt{spawn-sont} is a tuple of \texttt{\{txid\ vout\ offset\}}:

\begin{itemize}[noitemsep,topsep=2pt]

\item
  \texttt{txid}: ID of a Bitcoin transaction whose outputs contain
  the relevant \textsc{utxo}
\item
  \texttt{vout}: the relevant \textsc{utxo}'s position in a 0-indexed array of
  the ``vector of outputs'' from this transaction.
\item
  \texttt{offset}: relative ordinal number of a sat in this \textsc{utxo}; the
  Groundwire ID is pinned to this sat, which we will follow through
  successive \textsc{utxo}s.
\end{itemize}

To facilitate comets attesting to their identities via urbtc, we
introduced Groundwire's notion of multiple \textsc{pki}s to Jael, along with vane
tasks for registering and suspending \textsc{pki}s. Each corresponds to the name
of a Gall agent permissioned to update information for one \textsc{pki} (pending
userspace permissions work currently underway). We additionally
introduced a ``Suite C'' cryptographic core utilized by Groundwire
comets, detailed in \textsc{uip}-0136 (\citeyear{UIP-0136}).

When a ship receives an introductory ``self-attestation'' packet from a
Suite C comet, the comet includes the \textsc{pki} data needed for verification.
Ames then passes this data to Jael and subscribes for a verdict on
whether the self-attestation is valid. Jael passes the data to the
corresponding agent for that \textsc{pki}, and upon receiving a positive update
sends a fact to Ames along the relevant subscription,
and Ames completes the handshake appropriately. From here, the userspace
agent can continue to watch the \textsc{pki} substrate for further key rotations
for that comet.

Due to the confidential nature of key rotations on urbtc, userspace
agents only know upon seeing a key rotation that that identity's old
keys are now out-of-date; they do not know what the new networking
keys are until receiving a new self-attestation packet from that
identity. The Urbit ship must, then, ignore all new packets from that
identity except for self-attestation packets. We introduce the
\texttt{\%snob} task to Ames (as a counterpart to \texttt{\%snub}) which
is used to put an identity into this limbo state.

In addition to your sponsor---which maintains contact with
your ship for relay and route rediscovery from peers if your node is
behind a WiFi router---we extended the routing hint options in
\texttt{\$point:jael} with a \texttt{\$fief}, an IPv4, IPv6, or \textsc{dns}
static transport address that may be used for route selection in the
kernel and which we attest in onchain routing state on Bitcoin. Any node
that has such an address can straightforwardly sponsor and route for
comets running on mobile or a home server, and optionally supply
value-add services, so we're inclined to grow this set early to enable a
robust network topology and incentivize development of QoS protocols.

urbtc is implemented across \href{https://github.com/gwbtc/urbit}{\texttt{gwbtc/urbit}}, \href{https://github.com/gwbtc/vere}{\texttt{gwbtc/vere}},
and \href{https://github.com/gwbtc/groundwire}{\texttt{gwbtc/groundwire}}. We aim to upstream our Arvo and Vere work
into the corresponding \href{https://github.com/urbit/urbit}{\texttt{urbit/urbit}} and \href{https://github.com/urbit/vere}{\texttt{urbit/vere}}
repos.

\subsection{Confidential comets}\label{confidential-comets}

Our initial tweak was designed for a version of the urbtc protocol that
used ordinal inscriptions \citep{Rodarmor2026Inscriptions}. Inscriptions
serialize data as chained \texttt{OP\_PUSH} opcodes of a \textsc{utxo}'s Taproot
script in a ``commit'' transaction, then a ``reveal'' transaction makes
that hashed onchain data legible as a plaintext \textsc{mime} response. This
required all Urbit nodes to index all Bitcoin blocks to find urbtc reveal
transactions and use those to maintain their local \textsc{pki} state.

In testing we found that the computational burden of this indexing was
trivial for Groundwire nodes on a commercial \textsc{vps} and the developers' own
laptops. The main issue was that nodes could take hours to catch up with
their peers, which we ameliorated with a stopgap trusted snapshot system
for onboarding. This helped testers by skipping an initial sync from our
chosen block 943,140 to the present, but we couldn't help them speed up
recovery from unplanned downtime because Urbit doesn't allow users to
overwrite \textsc{pki} state outwith Jael's \textsc{api}. This snapshot system was just a convenience, and
injecting a null value where the snapshot should have been would trigger
re-indexing of the chain from block 943,140.

However, this question of using Bitcoin as arbitrary data storage has
been a source of long-running debate in the Bitcoin community, citing a
concern for keeping the burden on Bitcoin node operators as low as
possible \citep{BIP0110}. Making all \textsc{pki}
state public onchain is also a privacy concern, even if the data is
apparently trivial.

For these reasons, we implemented a system of ``confidential comets''
by, for a start, simply omitting the ``reveal'' transaction for
attestations. Instead, we have comets reveal their \textsc{pki} history offchain
as part of the initial handshake with a new peer. As of 8K, the three
data commitments that comprise a Groundwire ID attestation have the
following shapes.

\subsubsection{Name commitment}\label{name-commitment}

The keypair underlying the comet's name commits via Schnorr tweak
(scalar addition) to \texttt{{[}\%urbtc\ spawn-sont{]}}, which is a \textsc{pki}
tag and a unique chainstate pointer binding one name to one sat. (No
other name may be bound to this sat.)

The \texttt{\%urbtc} tag is a Hoon text constant that names our \textsc{pki} in
Urbit's Jael kernel module. The sat is addressed relative to its host
\textsc{utxo} at comet generation time rather than by a global sequence number as
in the ordinals protocol, due to heavy indexing requirements, but we
retain ordinals' first-in-first-out convention for ordering input sats
over transaction outputs when following a bound sat through its
subsequent spends.

\subsubsection{Onchain attestation}\label{onchain-attestation}

Subsequent \textsc{utxo}s containing the bound sat are used to permanently attest
the comet name and various \textsc{pki} state commitments beginning with the
spawn transaction, which is a spend of the original host \textsc{utxo} which we
create following key generation and sat selection.

This bidirectional cryptographic binding between the offchain (name) and
onchain (wallet) signers, plus the tweak's \textsc{pki} separator prefix,
prevents ``double-spending'' of a spawn transaction across \textsc{pki}s, as well
as some potential replay attacks within Bitcoin.

The base onchain ``confidential attestation'' is itself also a Schnorr
tweak, in Pay-to-Taproot (\textsc{p2tr}) output pubkeys. The scalar we add to the
wallet keypair to derive the \textsc{p2tr} key is a Merkle root committing to the
serialized Nock noun containing the comet name, its life (key rotation)
and rift (continuity-breaking key rotation), latest networking pubkey,
routing hints, and at least one of a sponsor (Urbit ID) or a fief
(\textsc{ip/dns}).

For an output hosting one ID and bound sat, its \textsc{pki} state is a Merkle
leaf at axis 2 (first left child), and currently this is the only
attestation type in the protocol. Batch custodial outputs containing
multiple bound sats (``shared \textsc{utxo}s'') aren't yet supported on the alpha
network, but we're exploring the design space particularly with respect
to potential L2 architectures. So except in case of a breach (Urbit's continuity-breaking key rotation)
the state hash is the root hash which is the scalar tweak. We reserve
the right-hand subtree for future protocol expansion to support other
attestation use-cases, such as display names and profile images,
external links, web-of-trust links, code signatures, and proof of funds.

At spawn the initial networking key is the (preimage of the) name
itself, but subsequent state commitments include both the name and a
distinct latest messaging pubkey. Every time an output bearing the sat
is spent, we require the new \textsc{utxo} containing it to commit the current
tip of its \textsc{pki} log, even if the spend is simple custody-management
e.\,g.\ transferring to a fresh address, and the committed state duplicates
the previous tip. Otherwise, the vocabulary of update types that we
support is 1) rotate networking key 2) update routing info 3) publish an
attestation and 4) continuity breach. Continuity breaches are
signalled by the presence of a 0-valued second leaf at axis 3.

Plaintext publication of \textsc{pki} data is an opt-in second layer ``public
attestation'' in an added \texttt{OP\_RETURN} output from the published
state transaction. This serves two purposes: 1) permissionless
advertising and discovery of public bootstrap sponsors, and 2)
peer-recovery-of-last-resort for ships partitioned from each other by a
mutual routing failure (e.\,g.\ extended sponsor downtime).

Note that we did not specify any tapscript encoding (Bitcoin Script in
Merkle leaves that a \textsc{p2tr} spend can execute with an inclusion proof) of
the protocol's Merkle data above. That would be necessary if we wanted
to publish data in the transaction witness to take advantage of the
Segwit fee discount, as in the ordinals inscription mechanism. However,
witness inscriptions bypass a Bitcoin node operator's ability to filter
the data without foregoing witness verification of signature operations
as well, and the Segwit discount is actively harmful in the Groundwire
setting since our sybil resistance mechanism \emph{is}
Bitcoin-denominated costliness.

Non-tapscript Merkle data commitments in the output chain are, instead,
verifiably revealed by an \texttt{OP\_RETURN} output that includes the
internal (pre-tweak) keys of previous outputs and their signatures along
with the state history. This is the same format as our offchain
revelation in the peer attestation packet, detailed below.

In the taxonomy of Bitcoin Layer 2 constructs, urbtc is therefore a
client-side validated (\textsc{csv}) protocol.

\subsubsection{Offchain revelation}\label{offchain-revelation}

In Urbit's peer-to-peer networking flow, comets introduce themselves offchain by exchanging attestation packets
with a peer, and provide a refreshed packet when they broadcast new
confidential attestations of \textsc{pki} state onchain.

Legacy comets simply sign a minimal handshake, sufficient to negotiate
the per-peer \textsc{ecdh} channel, using the same key of which their ID is a
straightforward hash. Groundwire IDs instead sign the packet with their
latest messaging key, \emph{and} must reveal the plaintext tweak data
their name committed to, plus all data required for the recipient's \textsc{pki}
agent to verify the signature, the validity of the key, and the validity
of the name. If the \textsc{pki} tag is recognized by the receiver's Jael, it
will forward the packet to the responsible \textsc{pki} agent and chain client
for validation. Otherwise, it will treat the sender as a legacy comet
and require it to be at life \texttt{1} and sign with its identity key. For urbtc
comets, the required information is outpoints and block heights for the
sender's entire onchain history, and plaintext state commitments with
Schnorr inclusion proofs for their current rift.

By the time we get to a production launch, ships should be able to
configure their own sybil resistance policies including minimum
acceptable values for costliness or proof-of-funds needed to accept a
peer. That any ship may have their own secret policy obscures the
requirements that a bad actor needs to meet to spawn a usable batch of
fake accounts for any one ship or range of ships under one policy. Since
a negative policy decision prevents peer discovery from happening at
all, ships may only ever further constrain whichever policy is applied
by their sponsor. Sponsors might advertise these policies as a point of
differentiation. Sponsees may route around their sponsors' policies by
accepting a new peer's keyfile out-of-band, e.\,g.\ by scanning a QR code
or tapping an \textsc{nfc} tag in-person.

Since the attestation packet format specifies the content, location, and
encoding of confidential attestations, and is itself broadcast in a
public attestation, we can use a single kelvin version on the packet to
cover both layers of the protocol, while freezing the ID tweak-data
format straight to 0K. The attestation packet's current kelvin version
is 8K.

\section[Wire \textsc{uri}s]{Wire URIs}\label{wire-uris}

The Wire \textsc{uri} scheme enables Groundwire IDs to make cryptographic claims
on other namespaces with minimal integration between Groundwire and that
namespace's ecosystem.

There are two forms of Wire \textsc{uri} (or ``wire''): signed, and unsigned.

\subsection{Unsigned wires}\label{unsigned-wires}

An unsigned wire is of the form:

\begin{lstlisting}[style=listingcode]
wire://id/tag/path
\end{lstlisting}

\noindent
as in these examples:

\begin{lstlisting}[style=listingcode]
wire://aback.baboon.cabal.enough/https/groundwire.io
wire://abyss.conceal.decrypt/geo/38.89278,-77.04583
\end{lstlisting}

An unsigned wire's \texttt{id} is a Groundwire ID in the form of a bare
nym, since we can't assume the signer and verifier exist on the same
\textsc{pki}. The \texttt{tag} (a string \texttt{\^{}{[}a-z{]}{[}a-z0-9-{]}*\$})
refers to some protocol that a client may or may not know what to do
with, and the \texttt{path} is any \texttt{/}-separated path that
corresponds to a resource the client can locate via that protocol.

An unsigned wire may be generated by the client as a human-readable reference to a
signed wire stored elsewhere in the client. Unlike signed wires, they do
not self-describe a version number, so a client parsing an unsigned wire
without reference to a corresponding signed wire would have to handle a
case if the \texttt{/tag/path} format ever changed during a kelvin version decrement in the
signed wire.

\subsection{Signed wires}\label{signed-wires}

Signed wires are of the form:

\begin{lstlisting}[style=listingcode]
wire://id/base64url
\end{lstlisting}

\noindent
as in this example:

\begin{lstlisting}[style=listingcode]
wire://abate.banal.campaign.debate.elite/8BEBAC9odHRw
       cy9leGFtcGxlLmNvbRam2y5UYO8BaCbZDGFToaaToJugs8
       BX_uEOs3Hy7M95LhTn21s3_Mg-wH9WG-CsUxm2JGWjAbou
       m7waoMSWfgc
\end{lstlisting}

Signed wires serve a similar function to \textsc{json} Web Tokens \citep{RFC7519}, in that we can
use them to non-interactively verify that the owner of some public key
signed a given string. We use Wire \textsc{uri}s and not \textsc{jwt}s in part because
Wire \textsc{uri}s have a human-readable trust root and always represent a path
in a signed namespace beneath that root.

Signed wires only have one segment after the \texttt{id}: a Base64URL
string that encodes a header and a payload.

The 12-bit header has this structure:

\begin{itemize}[noitemsep,topsep=2pt]

\item
  \textbf{First 1 bit:} A \texttt{0} or \texttt{1} flag that tells us if
  the signer has signed a hash of the \texttt{/tag/path} string
  (\texttt{1}), or the \texttt{/tag/path} string concatenated with the
  response we will receive from that path (\texttt{0}). This flag bit is
  not subject to kelvin versioning.
\item
  \textbf{Next 3 bits:} kelvin version number. Wires launch at 7K, or
  \texttt{111}. Everything to the right of this is version-dependent.
\item
  \textbf{Next 8 bits:} The bytewidth of the decoded \texttt{/tag/path}
  string, which is up to 256 \textsc{ascii} characters.
\end{itemize}

The payload has this structure:

\begin{itemize}[noitemsep,topsep=2pt]

\item
  \textbf{First 16 bits:} Key rotation number, an abstraction over
  whatever key rotation mechanism the \textsc{pki} substrate uses. May not be
  \texttt{0}, we assume the signer's first keypair is \texttt{1}.
\item
  \textbf{Next \emph{n} bytes:} The plaintext \texttt{/tag/path} string,
  where \emph{n} is specified in the header
\item
  \textbf{Next 64 bytes:} 7K's EdDSA signature of a 32-byte \textsc{sha}-256 of
  either \lstinline[style=inlinecode]{key_rotation || tag_path} or
  \lstinline[style=inlinecode]{key_rotation || tag_path || content_hash},
  where \lstinline[style=inlinecode]{content_hash} is itself a 32-byte hash of the octet
  stream a request on this path will receive. What exactly goes in the
  octet string depends on the \texttt{/tag} protocol; the signer may
  assume that a relevant client knows what kind of response to expect
  for a given protocol.
\end{itemize}

Clients can verify a signed wire that has no content on receipt. If they
receive a signed wire with content, they should reconstruct the content
hash on receipt and verify that against the signature before further
processing.

Since Groundwire ID persists across key rotations, the client must
verify the signature using the public key that corresponded to the
wire's \texttt{id} at the key rotation signed in the payload. Key rotation
events should not invalidate signed wires from previous key rotations,
as verifiers should be able to get the keypair that corresponds to the
wire from historical state. That historical state may be public
(e.\,g.\ onchain) or privately distributed among peers, relays, or other
network providers. (\textsc{pki}s such as urbtc may implement a special,
continuity-breaking key rotation event that \emph{does} invalidate
previous signatures, upon which peers should treat all historical state
for that ID as having been wiped. If an ID signed a wire at keypair
\texttt{2} before and after the continuity breach, only the
post-continuity-break keypair \texttt{2} should be under consideration
by the verifier. All signatures from before the breach should be considered invalid.)

Some example use-cases for Wire \textsc{uri}s are as follows:

\begin{itemize}[noitemsep,topsep=2pt]

\item
  \texttt{/mailto}: A business or government agency signs their email
  address in a public link.
\item
  \texttt{/https}: A business signs links in their emails such that a
  Groundwire-aware email client can filter out phishing scams without
  reference to a trusted third party.
\item
  \texttt{/tel}: A government agency signs their public-facing telephone
  numbers; they sign numbers they may call from such that a
  Groundwire-aware phone/VoIP client can verify the number is
  legitimate before it rings.
\item
  \texttt{/bitcoin}: A Groundwire ID requests payment to a Bitcoin
  address that may or may not hold the Groundwire ID's corresponding
  sat, enabling multi-wallet setups. The \textsc{uri} may also include a
  requested payment amount.
\item
  \texttt{/wss}: A WebSockets server provider links a server they run to
  their Groundwire ID, without necessarily having to run that server on
  their Groundwire node.
\item
  \texttt{/ssh}: A \textsc{vps} host signs their host key fingerprint,
  supplementing trust-on-first-use for users who want to \textsc{ssh} into that
  server.
\item
  \texttt{/geo}: A Groundwire ID signs coordinates of a geocache.
\item
  \texttt{/ipfs}: A Groundwire ID signs some distributed content by its
  hash, linking the content to the ID without the ID user having to run
  their own storage.
\end{itemize}

At time of writing there is no canonical reference implementation of the
Wire \textsc{uri} scheme, but it is used in the \texttt{gwbtc/chorus} repo to
store references to advertised content. This may be replaced with a
Kademlia \textsc{dht} implementation which stores references as content hashes
rather than Wire \textsc{uri}s, mostly to improve upon Chorus's gossip protocol.

\section{Future Work}\label{future-work}

Sections \S\ref{mnemonyms}, \S\ref{groundwire-id}, and \S\ref{wire-uris} describe code which is running in a public alpha
release. This section discusses three future projects which are most
pertinent to the ideas in this article: named data networking between
Groundwire IDs, post-quantum signatures, and ``child'' identities
derived from root-level Groundwire IDs.

\subsection{Named data networking}\label{named-data-networking}

Alongside kelvin versioning \S\ref{ossification}, the other part of Urbit that we aim
to lift out of that system is Directed Messaging (colloquially called
``Mesa'' by Urbit developers), which was detailed in a previous issue of
this journal \citep{DirectedMessaging2026}.

We think Mesa's named data networking (\textsc{ndn}) behavior is a natural fit
for any network of cryptographically sovereign root nodes, whether those
correspond to separate physical machines in a peer-to-peer network or as
virtual nodes on a cloud server.

\textsc{ndn} is sometimes described with terms like ``data-oriented'' as opposed
to ``connection-oriented'' \textsc{tcp/ip}, and one can think of Groundwire as a
data-oriented sybil resistance as opposed to connection-oriented sybil
resistance (\textsc{captcha}, JS challenge, etc.).

To pull Mesa out of Urbit means removing some of the assumptions around
network topology. For instance, Groundwire has no concept of a hierarchy
of ``galaxy'', ``star'', and ``planet'' nodes; sponsor-sponsee
relationships are formed between service providers and their users, and
any comet can sponsor another.

We don't assume that Groundwire nodes are running Urbit OS, though our
Mesa implementation should work for those nodes. One of the main goals
of our Mesa project is seamless peer-to-peer networking between
non-Urbit clients and Urbit servers.

Other than the caveats above and some other details, all headline
features of Directed Messaging should be implemented in a Mesa library
as protocol invariants:

\begin{itemize}[noitemsep,topsep=2pt]

\item
  All messages should be requests or authenticated responses carrying
  named data.
\item
  Authentication should happen on packets and not sessions.
\item
  Source-independent single hop routing should be sufficient for
  networking, which minimizes bookkeeping required for nodes.
\end{itemize}

Other features of Mesa, such as exactly-once messaging, may be
configurable per networking flow depending on the stack the peer in that
flow is running Mesa on. Details like cache management and congestion
control are implementation-specific.

This has the benefits of \textsc{ndn} as typically described, such as more
resilient routing in unpredictable physical network conditions and less
reliance on single points of control. Directed Messaging's
referentially-transparent namespace, with Groundwire IDs at the root,
has its own use-cases for the distributed networking cases we're
interested in:

\begin{itemize}[noitemsep,topsep=2pt]

\item
  Every path is signed by its publisher and has a version number, and
  there is currently no way to request data at a path without knowing
  its version number, so unversioned imports are impossible.
\item
  Referentially-transparent paths under identities with key rotation
  enable more long-lived authorities over links, and Mesa's caching and
  relay properties trivialize distributed archives and other \textsc{cdn}-like
  setups.
\item
  Cheap, mechanically trivial \textsc{cdn}s distribute load away from
  content publishers. Effectful writes must be signed by the
  Groundwire ID that originally sent them via \textsc{hmac}. The likelihood that
  a write or access-controlled read comes from a disposable ID in a botnet can be priced at the
  moment of receipt of the first packet: either the peer writing to us
  is known and validated, or we ask for a self-attestation packet that
  includes the \textsc{pki} evidence described in \S\ref{pki-evidence}.
\item
  Publishers and relays don't need to verify connections because there
  are no connections. Reads to access-controlled paths have similar
  caching properties as anonymous reads on public paths.
\end{itemize}

Our Mesa release will be a formal spec and at least one open-source
reference implementation, likely in C. Many implementation details are
left as exercises for network developers, including:

\begin{itemize}[noitemsep,topsep=2pt]

\item
  Peer discovery mechanisms (bootstrap nodes, \textsc{dht}s, pre-distributed
  \textsc{pki}s, etc.).
\item
  Transport layers (Mesa over \textsc{tcp}, \textsc{udp}, \textsc{sctp}, etc.).
\item
  Caching policies.
\item
  Peer-to-peer economic incentives (congestion control, relays, \textsc{cdn}s, etc.).
\item
  Interaction with anonymous requesters.
\item
  Reputation/uptime scoring.
\end{itemize}

Our Mesa work is in a preliminary \href{https://github.com/gwbtc/libmesa}{\texttt{gwbtc/libmesa}} repo, which
will be made public once it's ready for a public alpha release.

\subsection{Post-quantum signatures}\label{post-quantum-signatures}

At this point in time Groundwire's persistent cryptographic identity
across key rotations is aimed at networks like Nostr, which have no
affordance for transferring identity in the event of a compromised
private key. Key rotation is presently a narrow use-case
that even diligent users should rarely have to go through.

However, cryptographically-relevant quantum computing could make key rotation a much
more salient issue, even in normal circumstances. (PQ keypairs in
stateful schemes like \textsc{xmss} and \textsc{lms} are finite, enumerated trees of
single-use keypairs; the number of keypairs in a tree may be as low as
2\textsuperscript{5} or up to 2\textsuperscript{25} \citep{NIST800208}.) With Google recommending
migration to post-quantum cryptography by 2029 \citep{GoogleQuantumAI2026}, technical
communities like Bitcoin are urgently weighing their options for
migration of legacy wallets to \textsc{pqc}.

We use ``pre-quantum'' signatures for urbtc and Wire \textsc{uri}s at their
current kelvin versions because this is the status quo in both Urbit
networking and Bitcoin consensus. In particular, the Schnorr
cryptography we base our overlay of the comet namespace on is high among
the signature features that will be difficult or impossible to achieve
in the known PQ algorithms that are feasible for Bitcoin.

Our strategy for an early, low-overhead migration precisely mirrors the
\textsc{bip}-360 ``Pay-to-Merkle-Root'' quantum resistance project \citep{BIP0360}:

\begin{itemize}[noitemsep,topsep=2pt]

\item
  Abandon offchain key tweaking.
\item
  Breach comet semantics from pubkey hashes to pure hashes of structured
  data that attest multiple signature algorithms.
\item
  Implement a range of high-attention open PQ signature schemes drawn
  from Bitcoin research and other networks.
\end{itemize}

We expect that efforts in that third area, implementing new signature
schemes, are likely to lead and drive timing for the migration as a
whole, but a major goal will be to improve modularity of cryptosystem
adoption in the future as well as upgrading the core crypto suites.

Once Bitcoin developers and hardware providers settle on a PQ migration
plan, we will develop a migration plan from pre- to post-quantum kelvin
versions which preserve continuity of identity and attestations to the
fullest extent possible.

\subsection{Hierarchical identity}\label{hierarchical-identity}

Groundwire ID's namespace has no outer limit: a nym can have as many
words as you need. urbtc's namespace can accommodate \(2^{128}\)
identities because each must correlate to an Urbit comet. urbtc
identities must also correlate to Bitcoin sats, of which there will only
ever be 2.1 quadrillion in existence, with far fewer in circulation. How
could one Groundwire \textsc{pki} like urbtc scale to name all humans, devices,
and AI agents on and near Earth?

As it stands a Groundwire ID is an Ed25519 keypair generated by an HD
wallet, which gives us an obvious way to correlate ``sibling''
identities owned by the same user wallet. But it also allows urbtc to
fill the remainder of the 128-bit namespace by allowing the 2.1
quadrillion onchain root nodes to continue the next level of their HD
derivation tree and derive many ``child'' identities, each of which can
likewise derive their own children, and so on.

This hierarchical recursion of the HD wallet tree-structure gives us a
natural way to represent relationships of ownership, responsibility, and
granular access control amongst human-driven identities (e.\,g.\ multiple
signing devices belonging to one human) and AI agent identities
(e.\,g.\ multiple agents belonging to one human, multiple subagents
belonging to one agent, etc.). Since child keys are derived by tweaking
the parent with an arbitrary nonce, we can use this to commit to
descriptors for the role, origin, authority, etc. of the child,
cryptographic authorizations like a signature over permissions, or even
a discretionary spending budget in the form of Chaumian ecash
signatures.

We haven't explored the design space for this yet, but for our purposes
a system for deriving and verifying identities must satisfy these
requirements:

\begin{enumerate}[noitemsep,topsep=2pt]
\def\labelenumi{\arabic{enumi}.}

\item
  Derived identities are trivially verifiable as belonging to their
  ultimate root identities and any parent identities in between.
\item
  Deriving an identity is mechanically simple, computationally cheap,
  and almost instantaneous; an AI agent can trivially derive their own
  identities, having gained approval from its root identity via manual
  sign-off or hard-coded policy.
\item
  Freshly derived identities have no sybil resistance value of their
  own, instead backed by the sybil value of their ultimate root identity
  and subject to whichever calculation a verifier chooses.
\item
  Parent/child relationships can be expressed in the namespace and
  children may be identified via their parents; e.\,g.\ if AI agent
  \texttt{..denounce...escape} belongs to human operator
  \texttt{.abet...baboon}, they might be identified in a UI as
  \texttt{.abet...baboon::denounce...escape}, or
  \texttt{::abet...escape}, or some other representation. Namespaces
  belonging to a derived identity may be viewed as subtrees of their
  root identity's namespace.
\end{enumerate}

\section{Conclusion}\label{conclusion}

Groundwire is cool and you should check it out.\tombstone

\section{Appendices}\label{appendices}

The following appendices are supplementary notes on software
ossification and consensus mechanisms, detailing the thinking behind
some of our design decisions above.

\subsection{Ossification}\label{ossification}

One of the premises of open-source software is that ``given enough
eyeballs, all bugs are shallow''. Looking at recent large-scale
cyberattacks---like the 2025 ``Shai-Hulud'' attacks distributed via
node package manager (npm), and possibly written by an \textsc{llm} \citep{Moore2025ShaiHulud}---we anticipate
a future where the majority of those eyes are fleets of fully-automated
cyberattack agents with near- or super-human ability to discover
vulnerabilities and exploit targets via malicious code and social
engineering. All automatic software update pipelines are attack vectors,
and dependency managers such as npm might import dozens of untrusted
dependency versions---recursively---at build time for even the
smallest projects.

Conversely, the developers of these dependencies find that updates
become more fraught the greater the number of dependents, such that even
desirable updates become impossible in practice. This kind of unplanned
ossification is almost always bad. Two software communities in which
developers embrace planned ossification are Urbit, which formally counts
down to an immutable OS kernel, and Bitcoin, where an informally-defined
stable state is seen as part of the investment thesis.\footnote{Jameson Lopp
argued against an ossifying Bitcoin in 2024 \citep{Lopp2024}, but gave a
generous summary of the pro-ossification arguments he'd encountered.}

In a 2017 blog post on \texttt{urbit.org} called ``Towards a Frozen Operating
System'' \citep{YarvinWolfePauly2017}, Curtis Yarvin
and Galen Wolfe-Pauly contended that ``there is almost never any good
reason to automatically update dependencies with fresh versions of
foreign code. In the case of a security update, the maintainer of an app
which depends on a patched library needs to catch the patch and
propagate the update in a new version of the app. An app shouldn't be
security-critical to begin with, and propagating updates is a basic
function of a modern OS. Also, if the app is security-critical, it could
have its own bugs, so it needs a maintainer and an update pathway
anyway.'' In a 2020 \texttt{urbit.org} post \citep{Monk2020} Urbit developer
Philip Monk wrote that ``system modules should be designed to be frozen
permanently; conversely, Urbit consists of that system software which
admits of a Kelvin version.''

What is a kelvin version? The 2016 Urbit whitepaper \citep{Whitepaper} doesn't define
it. At time of writing, no formal specification has been signed off on
by the Urbit Foundation or Urbit's core developers. We do find one
subsubsection of informal description in Yarvin's 2010 blog post
\emph{Urbit: functional programming from scratch}, which publicized
Urbit under the name ``Urbit'' for the first time: ``the true, Martian
way to perma-freeze a system is what I call Kelvin versioning. In Kelvin
versioning, releases count down by integer degrees Kelvin. At absolute
zero, the system can no longer be changed. At 1K, one more modification
is possible. And so on.'' \citep{Yarvin2010}

Today, Urbit's OS is distributed at 408K. This version number tags the
\texttt{zuse.hoon} file which defines Urbit's kernel module \textsc{api}s; the
modules themselves are not kelvin versioned. Zuse is the ``outermost''
layer of Urbit's OS, and deeper layers may have their own kelvin
versions. (The Hoon standard library at 135K, for instance.)

This raises the question of kelvin-versioned dependency management. In
``Towards a Frozen Operating System,'' Yarvin and Pauly wrote:

\begin{quote}
The right way for this trunk to approach absolute zero is
to ``telescope'' its Kelvin versions. The rules of telescoping are
simple:

\begin{enumerate}[noitemsep,topsep=2pt]
\item
  If tool B sits on platform A, either both A and B must be at absolute
  zero, or B must be warmer than A.
\item
  Whenever the temperature of A (the platform) declines, the
  temperature of B (the tool) must also decline.
\item
  B must state the version of A it was developed against.
\item
  A, when loading B, must state its own current version, and the
  warmest version of itself with which it's backward-compatible.
\item
  Of course, if B itself is a platform on which some higher-level tool
  C depends, it must follow the same constraints recursively.
\end{enumerate}

The important effect of obeying rule 2: in a tall trunk,
it's sufficient to state your dependency on the platform directly below
you. Suppose C sits on B, which sits on A and exports features of A. C
can state only the version of B it's written against, because when A
chills, B must also chill.
\end{quote}

In practice, Urbit core developers update Zuse and its dependencies at
once: every Urbit release is either a new Zuse kelvin or a non-kelvin patch to
the current distribution of Zuse. Any change to any kelvin-versioned
file below Zuse will be slated for release in the next Zuse kelvin release.
Significant changes to unversioned kernel modules usually go in the next
Zuse release, but this is at the discretion of the Urbit Core Guild.

Today, Urbit userspace applications must declare the Zuse version they
were built against. If the user is running Zuse 408K, an application
update built against 407K will be received and stored for later. Once
the user is ready to update to Zuse 407K, the app update is applied in
the latter stages of an \emph{atomic} upgrade of Zuse and its
dependents, which will cleanly fail if even one dependent has no pending
update to 407K. If the user wishes to update Zuse anyway, they can
``suspend'' the application to prevent it from running but hang onto its
code and data in the event that a 407K update ever arrives.

The maintenance burden this places on Urbit userspace developers is
nontrivial: most Urbit userspace code available online will not compile
against Zuse 408K. This fact runs counter to one of the stated goals of
kelvin versioning, which is perfect forward compatibility. Quoting
``Towards a Frozen Operating System'':

\begin{quote}
C is a very mature language. C has long since mastered
backward compatibility. We can be sure that any C program which compiles
on today's correct C compilers will compile on any future correct C
compiler.

But, since C is not frozen, we can't be sure that today's
C compilers will compile any future C program. So C has not yet achieved
perfect forward compatibility. By definition, a frozen system is
compatible both forward and backward. So C is mature, but it could still
get more mature.

The point of a frozen OS is to shoot for infinite
maturity. Obviously a frozen platform is compatible forever in both
directions.
\end{quote}

The Urbit Foundation is doing prototype work on ``kelvin shimming'',
which allows Urbit to run userspace applications built against previous
kelvins in their own sandbox, but what form that solution will take is
still an open question.

\subsection{Consensus}\label{consensus}

Groundwire is a protocol and has no opinions about technical substrates.
The Groundwire Foundation has such opinions, and our choice of the
Bitcoin blockchain as the substrate for urbtc is highly motivated by the
same concerns that inform our protocol design.

The most important feature of Bitcoin for our use-case is the
``longest-chain rule'' in its consensus protocol. Per the Bitcoin
whitepaper \citep{Nakamoto2008}, ``proof-of-work
\ldots{} solves the problem of determining representation in majority
decision making. If the majority were based on one-IP-address-per-vote,
it could be subverted by anyone able to allocate many IPs. Proof-of-work
is essentially one-\textsc{cpu}-one-vote. The majority decision is represented by
the longest chain, which has the greatest proof-of-work effort invested
in it.''

A Bitcoin node accepts the longest proof-of-work chain as the canonical
record of Bitcoin's transaction history, without question. If one
finalized chain is overwritten by a longer one, so be it. This simple
decentralized governance mechanism has survived many bitter disputes in
Bitcoin's history. Amongst competing Bitcoin forks, a single chain holds
onto almost all nodes due to simple network effects and the
game-theory of proof-of-work: in 2014 Vitalik Buterin wrote that ``when
contributing one's work to the blockchain, a miner must make the choice
of which of all possible forks to contribute to (or whether to try to
start a new fork), and the different options are mutually exclusive.
Double-voting, including double-voting where the second vote is made
many years after the first, is unprofitablem {[}sic{]} since it requires
you to split your mining power among the different votes; the dominant
strategy is always to put your mining power exclusively on the fork that
you think is most likely to win.'' \citep{Buterin2014}

Technical debates amongst Bitcoin developers, such as that around \textsc{bip}-110 (\citeyear{BIP0110}), hold it as a
priority that independent and hobbyist node operators using consumer
hardware should be able to feasibly replay the history of the chain to
arrive at its current state. This property is essential for
decentralization and it falls out of proof-of-work: on initial sync a
node operator may choose to trust or verify previously-committed blocks
via an \texttt{assumevalid} flag, but in any case they are replaying the
longest proof-of-work chain upon the genesis block which serves as the
trust root.

Quoting \textsc{uip}-0136, ``Groundwire's thesis is pragmatic: over the long
term, Bitcoin is the only chain that provides an unambiguous,
protocol-level canonical history (the heaviest valid proof-of-work
chain) that a Kelvin 0 Urbit could resolve without social coordination
or trusted checkpoints. This suggests a clean long-term anchor for the
Urbit \textsc{pki}.''

Compare this to Ethereum's proof-of-stake model. If an Ethereum node
receives two finalized blocks that contradict each other, there is no
programmatic fork-choice rule to choose between them \citep{EthereumOrgWeakSubjectivity}. The node's
version of chain history must be re-derived from a ``weak subjectivity
checkpoint'': some block height acquired from a trusted third party
which they claim to be a truth universally acknowledged by all node
operators.

\textsc{uip}-0136 continues: ``It is too early to propose that that belief be
enforced by {[}Urbit{]} kernel policy \ldots{} Thus while Bitcoin's
popular ordinal/inscription protocol is a functional and
easy-to-reason-about initial choice of onchain mechanism, it would also
be premature to commit the kernel to these abstractions. Instead, Urbit
should provide a general mechanism whereby comets can attest keys to any
\textsc{pki} source (chosen by a Jael \%listen task) and peers can verify these
attestations.''

Having enabled Groundwire users to compose multiple \textsc{pki} substrates---at
least one of which is Bitcoin, and several of which may be other
blockchains---we have fallen through a trapdoor and landed in the
``blockchain interoperability problem''. Groundwire ID's name commitment
is supposed to namespace \textsc{pki}s such that two \textsc{pki}s should never hold the
same ID, but if we nonetheless receive a self-attestation packet from an
ID on Solana claiming to be an ID on Bitcoin we are about to send money
to, who do we trust? Which of our \textsc{pki} gateways is broken or malicious?
Which record of history is correct? What are the semantics by which
history is recorded, and how do we reconcile those histories if two
consensus mechanisms use different semantics?

At this point we may fall back on human judgement, but Groundwire
assumes that 99.9\% of internet users are not humans and cannot be held
accountable for their judgements.

If we grant that even Bitcoin's longest-chain rule is inevitably backed
by the human drama of social consensus, institutional governance,
conflicting incentives, and spirited public debate, we think this is all
the more reason to secure those out-of-band techno-social systems
against a flood of highly capable bad actors with kelvin-versioned
software and sybil-resistant identity and networking. By hosting the \textsc{pki}
onchain, Bitcoin could secure the host environment in which its own
developers, node operators, and other stakeholders work in public and
come to important decisions about its future. This holds for
proof-of-stake chains which are even more dependent on these out-of-band
consensus mechanisms.

So a programmatic symmetry-breaking rule is desirable where we cannot
reconcile the histories of two blockchains. If the chains in question
are both Ethereum L2s, we wait for the transactions to settle on
Ethereum; if Ethereum produces two finalized blocks, we phone a friend.
The opening of the Bitcoin whitepaper gives us such a programmatic rule,
which we may apply recursively upwards through increasingly meta-level
conflicts amongst consensus mechanisms: ``we propose a solution to the
double-spending problem using a peer-to-peer network. The network
timestamps transactions by hashing them into an ongoing chain of
hash-based proof-of-work, forming a record that cannot be changed
without redoing the proof-of-work. The longest chain not only serves as
proof of the sequence of events witnessed, but proof that it came from
the largest pool of \textsc{cpu} power. As long as a majority of \textsc{cpu} power is
controlled by nodes that are not cooperating to attack the network,
they'll generate the longest chain and outpace attackers.''

May the longest proof-of-work chain win.

\selectlanguage{USenglish}
\printbibliography
\end{document}
